\section{Recovering the full Clifford correspondence at a general point}
\label{gen:section}

The argument in this section uses the equality
\begin{equation}\label{gen:full-group}
 \operatorname{Hg}(V)=\operatorname{SO}(V,q),
 \qquad V=h^\perp\subset H^2(S,\mathbb Q),\qquad \dim V=21.
\end{equation}
There is no genericity assumption on the abelian variety that supplies the
initial algebraic correspondence.  The principal point is to recover every
multiplicity in the full Clifford construction from one nonzero map.

\subsection{The generic Hodge group}

\begin{proposition}\label{gen:genericity}
In any sufficiently small, full-dimensional marked polarized period chart,
Equation~\eqref{gen:full-group} holds outside a countable union of proper
closed analytic subsets.  At every point satisfying
Equation~\eqref{gen:full-group}, the Hodge structure $V$ is irreducible and
$T(S)=V$.
\end{proposition}

\begin{proof}
Polarized local Torelli identifies such a chart with an open subset of the
domain of oriented positive two-planes in the fixed real quadratic space
$V_{\mathbb R}$; see \cite[Chapter~6]{Huybrechts}.
The infinitesimal generators of the Hodge circles over any such open
subset span $\mathfrak{so}(V_{\mathbb R},q)$.  Indeed, under the natural
identification $\mathfrak{so}(V_{\mathbb R},q)\simeq\Lambda^2V_{\mathbb R}$,
each generator is a nonzero multiple of the bivector of its positive
plane.  A linear functional vanishing on all these bivectors is an
alternating bilinear form vanishing on an open set of ordered pairs of
vectors.  It therefore vanishes identically.

The Hodge group is a connected rational algebraic subgroup of
$\operatorname{SO}(V,q)$.  Connectedness follows because its rational
identity component contains the connected Hodge circle and hence, by
minimality, is the whole group.  If this subgroup is proper, its Lie
algebra is a nonzero proper rational subspace
$L\subset\mathfrak{so}(V,q)$.  Set $r=\dim L$, and choose a rational
generator $\tau_L$ of $\Lambda^rL$.  Conjugation by the circle preserves
$L$ and acts trivially on its real determinant line: a continuous
character from a compact connected circle to $\mathbb R^\times$ is
trivial.  Thus $\tau_L$, viewed in
$\Lambda^r\operatorname{End}(V)$, has Hodge type $(0,0)$.

There are only countably many possible rational subspaces $L$.
For each one, the condition that $\tau_L$ have type $(0,0)$ is closed
analytic in the period chart.  In fact the tensor has weight zero and is
real, so this condition is equivalent to $\tau_L\in F^0$, a holomorphic
condition.  The locus is proper.  Otherwise $L_{\mathbb R}$ would be
preserved by every circle in the chart, hence by the adjoint actions of
their infinitesimal generators.  The spanning assertion above would
make $L_{\mathbb C}$ an ideal of the simple Lie algebra
$\mathfrak{so}_{21}(\mathbb C)$, a contradiction.

Finally, the standard representation of $\operatorname{SO}(V,q)$ is
absolutely irreducible and has no nonzero invariants.  Consequently $V$
has no proper nonzero rational Hodge substructure and no nonzero
rational vector of type $(1,1)$.  Since the primitive part of
$\operatorname{NS}(S)_{\mathbb Q}$ consists of such vectors,
$\operatorname{NS}(S)_{\mathbb Q}=\mathbb Qh$ and $T(S)=V$.
\end{proof}

\subsection{The spin representation forced by weight one}

Assume Equation~\eqref{gen:full-group}.  Let $B$ be a complex abelian
variety, put $U=H^1(B,\mathbb Q)$, and write
\[
 W^+=C^+(V,q),\qquad W^-=C^-(V,q)
\]
for the full Clifford parity spaces with their weight-one Hodge
structures.  The Hodge circle on these spaces acts by left multiplication
by its lift to $\operatorname{Spin}(V,q)$.
Let $\mathcal H$ be the joint Hodge group of $V,U,W^+,W^-$, and let
$\mathfrak h$ be its rational Lie algebra.

\begin{lemma}\label{gen:joint-group}
There is a decomposition into commuting rational ideals
\[
 \mathfrak h=\mathfrak g\oplus\mathfrak k,
 \qquad \mathfrak g\simeq\mathfrak{so}(V,q),
\]
where $\mathfrak k$ is the kernel of the action on $V$.
The ideal $\mathfrak k$ acts trivially on $W^+$ and $W^-$, and
$\mathfrak g$ acts on them by the usual left spin action.
\end{lemma}

\begin{proof}
The group $\mathcal H$ is connected by the same circle argument as in
Proposition~\ref{gen:genericity}.  It is reductive.  To see this directly,
each of the four pure summands is semisimple as a rational
$\mathcal H$-module: an invariant rational subspace is a Hodge
substructure, and a polarization provides a rational Hodge complement.
This is the usual semisimplicity argument for polarizable Hodge
structures; compare \cite[Proposition~3.6]{Deligne}.  Their direct sum is
a faithful semisimple module.  The unipotent radical has nonzero
invariants in each irreducible constituent; normality makes those
invariants $\mathcal H$-stable, so it acts trivially on every
constituent.  Faithfulness forces the radical to be trivial.

Projection to $V$ maps $\mathcal H$ onto
$\operatorname{Hg}(V)=\operatorname{SO}(V,q)$.
Indeed, its closed rational image contains the projected circle and
hence contains $\operatorname{Hg}(V)$. Conversely, the preimage of
$\operatorname{Hg}(V)$ contains the joint circle, so minimality places
$\mathcal H$ inside that preimage.
The reductive rational Lie algebra $\mathfrak h$ is
the direct sum of its center and its rational simple ideals.  Exactly
one simple ideal maps nontrivially to the absolutely simple quotient
$\mathfrak{so}(V,q)$: its image is the whole quotient, and the commuting
images of all other ideals must vanish.  The nonzero map from that
simple ideal is an isomorphism.  This gives the asserted rational
decomposition.

On $V\oplus W^+\oplus W^-$, the image of $\mathcal H$ is contained in
the rational algebraic image of $\operatorname{Spin}(V,q)$ acting on
$V$ through its orthogonal quotient and on $W^\pm$ by left
multiplication.  That image contains the projected Hodge circle, so
the containment again follows by minimality.  The differential of
$\operatorname{Spin}(V,q)\to\operatorname{SO}(V,q)$ is an isomorphism.
Its kernel on Lie algebras is therefore zero, proving both remaining
assertions.
\end{proof}

Let $\Delta$ denote the complex spin representation of
$\mathfrak{so}_{21}(\mathbb C)$, of dimension $2^{10}$.

\begin{lemma}\label{gen:width}
There are complex multiplicity spaces $M,N_+,N_-$ and a trivial
$\mathfrak g_{\mathbb C}$-module $M_0$ such that
\begin{equation}\label{gen:spin-decompositions}
 U_{\mathbb C}=(\Delta\otimes M)\oplus M_0,
 \qquad W^\pm_{\mathbb C}=\Delta\otimes N_\pm,
 \qquad \dim N_\pm=2^{10}.
\end{equation}
Every rational $\mathfrak g$-equivariant map $U\to W^\pm$ is a
weight-one Hodge map.  Consequently
\begin{equation}\label{gen:hom-base-change}
 \operatorname{Hom}_{\mathrm{HS}}(U,W^\pm)\otimes_{\mathbb Q}\mathbb C
 =\operatorname{Hom}_{\mathfrak g_{\mathbb C}}
       (U_{\mathbb C},W^\pm_{\mathbb C})
 \simeq\operatorname{Hom}_{\mathbb C}(M,N_\pm).
\end{equation}
\end{lemma}

\begin{proof}
Let $D$ be the infinitesimal Hodge-circle generator divided by $i$;
its eigenvalue on type $(p,r)$ is $p-r$.  In the decomposition of
Lemma~\ref{gen:joint-group}, write $D=D_{\mathfrak g}+D'$.  On
$V_{\mathbb C}$ the eigenvalues are $2,-2,0,\ldots,0$.
In the standard Cartan coordinates for type $B_{10}$, in which the
vector weights are $\pm e_1,\ldots,\pm e_{10},0$, we may therefore take
$D_{\mathfrak g}=2e_1^\vee$.

The element $D_{\mathfrak g}$ is semisimple on every finite-dimensional
$\mathfrak g_{\mathbb C}$-module.  It commutes with $D$, so
$D'=D-D_{\mathfrak g}$ is also semisimple on $U_{\mathbb C}$.
Each $D'$-eigenspace is $\mathfrak g_{\mathbb C}$-stable.
Consider an irreducible constituent in such an eigenspace, with
highest weight $\lambda=(\lambda_1,\ldots,\lambda_{10})$.
The highest-weight description for $B_{10}$ gives
\[
 \lambda_1\ge\cdots\ge\lambda_{10}\ge0,
 \qquad
 \lambda_i\in\mathbb Z\ \text{for every }i,
 \quad\text{or}\quad
 \lambda_i\in\tfrac12+\mathbb Z\ \text{for every }i.
\]
The largest and smallest eigenvalues of $2e_1^\vee$ are respectively
$2\lambda_1$ and $-2\lambda_1$.  Indeed $e_1^\vee$ is dominant, and the
Weyl group permits every coordinate sign change.  These elementary
highest-weight and spin facts may be found in the orthogonal-group
chapters of \cite{FultonHarris}.

Adding the scalar by which $D'$ acts does not change the difference
$4\lambda_1$ between the extreme eigenvalues.  Since $U$ has Hodge
types $(1,0)$ and $(0,1)$, the eigenvalues of $D$ on $U_{\mathbb C}$
belong to $\{1,-1\}$.  Hence $4\lambda_1\le2$.  The only possibilities
are the trivial representation and
$\lambda=(\tfrac12,\ldots,\tfrac12)$, which is $\Delta$.
On a spin constituent the extreme eigenvalues already are $1,-1$;
the only scalar shift preserving this pair is zero.  Thus
\begin{equation}\label{gen:kernel-circle-zero}
 D'=0\quad\text{on }\Delta\otimes M.
\end{equation}
This proves the assertion about $U$.  The complex Clifford identity
$C^+(V_{\mathbb C},q)\simeq\operatorname{End}_{\mathbb C}(\Delta)$,
with left multiplication, gives the assertion about $W^+$.
An invertible odd right multiplier identifies the two parity spaces
for the left spin action and gives the assertion about $W^-$.
Lemma~\ref{gen:joint-group} shows that $D'=0$ on both of them.

Now a rational $\mathfrak g$-map $f:U\to W^\pm$ vanishes on $M_0$
after complexification.  The summand $M_0$ is $D$-stable.  On the spin
summand, Equation~\eqref{gen:kernel-circle-zero} shows that $f$
intertwines $D$.  It therefore intertwines $D$ on all of $U$, hence
the Hodge circle, and is a Hodge map.  Conversely, every rational
Hodge map intertwines $\mathcal H$ and therefore $\mathfrak g$.

Finally, rational $\mathfrak g$-intertwiners are the kernel of the
finite system of rational linear equations
\[
 f\rho_U(X_a)=\rho_{W^\pm}(X_a)f,
\]
where $\{X_a\}$ is a rational basis of $\mathfrak g$.
Taking that kernel commutes with extension of scalars.  Schur's
Lemma over $\mathbb C$ and Equation~\eqref{gen:spin-decompositions}
now give Equation~\eqref{gen:hom-base-change}.
\end{proof}

\begin{lemma}\label{gen:vector-multiplicity}
The space
$\operatorname{Hom}_{\mathfrak g_{\mathbb C}}
 (V_{\mathbb C},\Delta\otimes\Delta)$ has dimension one.
Neither $\Delta$ nor the trivial representation contains a vector
copy $V_{\mathbb C}$.
\end{lemma}

\begin{proof}
The spin representation in type $B_{10}$ is self-dual, since the
longest Weyl element is $-1$.  Under the conjugation action, the
Clifford antisymmetrization map gives
\begin{equation}\label{gen:clifford-exterior}
 \Delta\otimes\Delta\simeq\operatorname{End}(\Delta)
 \simeq C^+(V_{\mathbb C},q)
 \simeq\bigoplus_{k\ \mathrm{even}}\Lambda^kV_{\mathbb C}.
\end{equation}
Here the action is conjugation, in contrast to the left action used
in Equation~\eqref{gen:spin-decompositions}.

For completeness, the vector multiplicity in
Equation~\eqref{gen:clifford-exterior} can be read off without
decomposing every exterior power.  Choose an orthonormal basis
$v_1,\ldots,v_{21}$.  An equivariant map
$F:V_{\mathbb C}\to\Lambda^kV_{\mathbb C}$ is determined by $F(v_1)$,
because the orthogonal orbit of $v_1$ spans $V_{\mathbb C}$.
Its value is fixed by all determinant-one diagonal sign changes on
$v_2,\ldots,v_{21}$.  A basis exterior monomial fixed by all these
sign changes contains either none or all of these twenty vectors.
For even $k$, only $k=0$ and $k=20$ can therefore occur.
The degree-zero map vanishes, since a determinant-one orthogonal
transformation sending $v_1$ to $-v_1$ forces
$F(v_1)=F(-v_1)=-F(v_1)$.  For $k=20$ the possible value is a multiple
of $v_2\wedge\cdots\wedge v_{21}$, and volume duality
$\Lambda^{20}V_{\mathbb C}\simeq V_{\mathbb C}$ supplies a nonzero
equivariant map.  The multiplicity is exactly one.
The last assertion follows from irreducibility and the distinct
dimensions $21$, $2^{10}$, and $1$.
\end{proof}

\subsection{Independent target operations and rational descent}

\begin{theorem}\label{gen:generation}
Assume Equation~\eqref{gen:full-group}.  Suppose that
$\alpha\in\operatorname{CH}^2(S\times B)_{\mathbb Q}$ induces a
nonzero map $V\to H^2(B,\mathbb Q)$ for some complex abelian variety
$B$.  Let $\epsilon_1,\epsilon_2\in\{+,-\}$, and let $A_1,A_2$ be
abelian varieties realizing $W^{\epsilon_1},W^{\epsilon_2}$ as their
first rational cohomology.  Every rational Hodge map
\[
 \beta:V\longrightarrow W^{\epsilon_1}\otimes W^{\epsilon_2}
 \subset H^2(A_1\times A_2,\mathbb Q)
\]
is the restriction of the action of a cycle
$\Gamma\in\operatorname{CH}^2(S\times A_1\times A_2)_{\mathbb Q}$.
In particular, this recovers the full primitive Clifford
multiplication tensor and its prescribed even-Clifford realization.
\end{theorem}

\begin{proof}
Pull back $\alpha$ along addition $B\times B\to B$ and apply the
algebraic projector to $U\otimes U\subset H^2(B\times B,\mathbb Q)$
from Lemma~\ref{cyc:projectors}.  The resulting codimension-two cycle
induces
\[
 \alpha':V\longrightarrow U\otimes U.
\]
It is nonzero: on $H^2(B,\mathbb Q)=\Lambda^2U$, the mixed component
of addition pullback is the injective map
$u\wedge v\mapsto u\otimes v-v\otimes u$.

Fix a nonzero $\mathfrak g_{\mathbb C}$-equivariant map
$\gamma:V_{\mathbb C}\to\Delta\otimes\Delta$.
The map $\alpha'$ is rational and Hodge, hence
$\mathfrak g$-equivariant.  Lemmas~\ref{gen:width} and
\ref{gen:vector-multiplicity} give
\begin{equation}\label{gen:initial-tensor}
 \alpha'_{\mathbb C}=\gamma\otimes t,
 \qquad 0\ne t\in M\otimes M.
\end{equation}
In this formula the tensor factors have been regrouped so that the
two spin factors come first.  The other summands of
$U_{\mathbb C}\otimes U_{\mathbb C}$ contain no vector copy.

Put $N_j=N_{\epsilon_j}$ for $j=1,2$.  The independent operations
\[
 a\in\operatorname{Hom}_{\mathbb C}(M,N_1),\qquad
 b\in\operatorname{Hom}_{\mathbb C}(M,N_2)
\]
send $t$ to $(a\otimes b)t$, and these images span
$N_1\otimes N_2$.  To check this, choose
$\mu,\nu\in M^\vee$ with $c=(\mu\otimes\nu)(t)\ne0$.
For arbitrary $n_1\in N_1$ and $n_2\in N_2$, the maps
\[
 a(m)=\mu(m)n_1,\qquad b(m)=\nu(m)n_2
\]
send $t$ to $c\,n_1\otimes n_2$.  Thus any symmetry imposed on $t$
by the original exterior square causes no restriction: the two
operations are chosen independently.

To implement these operations over $\mathbb Q$, choose rational bases
\[
 f_1,\ldots,f_r\quad\text{of }
 \operatorname{Hom}_{\mathfrak g}(U,W^{\epsilon_1}),\qquad
 g_1,\ldots,g_s\quad\text{of }
 \operatorname{Hom}_{\mathfrak g}(U,W^{\epsilon_2}).
\]
By Lemma~\ref{gen:width}, all these maps are Hodge, and their complex
spans are the full spaces of the independent operations just used.
Lemma~\ref{cyc:h1} realizes each by a rational abelian homomorphism
acting contravariantly on $H^1$; this is the classical weight-one
equivalence, also stated in \cite[Theorem~6.8 and Corollary~6.9]{Milne}.
Products of their transpose-graph correspondences, followed by the mixed
projector if necessary, therefore realize every map
\begin{equation}\label{gen:rational-generators}
 \Phi_{ij}=(f_i\otimes g_j)\alpha'
\end{equation}
by a rational algebraic cycle $\Gamma_{ij}$ on
$S\times A_1\times A_2$.
These cycles have codimension two: the target correspondences have
degree shift zero, and composition preserves the degree shift of
the original surface correspondence.

Lemma~\ref{gen:vector-multiplicity},
Equation~\eqref{gen:initial-tensor}, and the spanning calculation show
that the complex span of the maps in
Equation~\eqref{gen:rational-generators} is
\[
 \operatorname{Hom}_{\mathfrak g_{\mathbb C}}
 \bigl(V_{\mathbb C},
       W^{\epsilon_1}_{\mathbb C}\otimes
       W^{\epsilon_2}_{\mathbb C}\bigr).
\]
The prescribed rational Hodge map $\beta$ lies in this space.
Consequently the finite system
\[
 \beta=\sum_{i=1}^r\sum_{j=1}^s c_{ij}\Phi_{ij}
\]
has a solution over $\mathbb C$.  All its coefficients and its
right-hand side are rational, so it has a solution with
$c_{ij}\in\mathbb Q$ by Lemma~\ref{cyc:descent}.  The cycle
$\Gamma=\sum_{i,j}c_{ij}\Gamma_{ij}$ has precisely the required action.

Finally, Clifford multiplication is the rational Hodge map
\[
 V(1)\longrightarrow\operatorname{Hom}(W^+,W^-),
 \qquad x\longmapsto L_x.
\]
Using a weight-one polarization to identify
$(W^+)^\vee(-1)$ with $W^+$ makes it a rational Hodge map
$V\to W^-\otimes W^+$.  It is therefore one of the maps just realized.
The fixed odd right multiplier and prescribed weight-one duality
then give exactly the specified even-Clifford tensor, by
Lemma~\ref{cliff:slots}.  No choice of a single spin copy or
undetermined scalar remains.
\end{proof}

\begin{remark}\label{gen:rationality}
The complex decompositions in
Equation~\eqref{gen:spin-decompositions} are used only to compute a
span.  The maps $f_i,g_j$ and the cycles $\Gamma_{ij}$ are rational
before that computation.  In particular, the argument requires
neither a split rational Clifford algebra nor descent of an
individual complex spin summand.  It proves algebraicity of the
prescribed rational tensor on the full Clifford spaces.
\end{remark}
